\documentclass[10pt]{article}
\usepackage[T1]{fontenc}
\usepackage[utf8]{inputenc}
\usepackage[french]{babel}
\usepackage{amsmath,amssymb,mathrsfs,mathpazo}
\usepackage{geometry}
\geometry{a4paper,margin=2cm}
\usepackage{graphicx,array,tabularx,tabularray,colortbl,multirow,multicol,booktabs,enumitem,pifont,ulem,fancybox,diagbox,cancel}
\usepackage{pstricks,pstricks-add,pst-plot,pst-tree,pst-node,pst-text,pst-func,pst-math,pst-bspline,pst-3dplot,pst-eucl,pst-solides3d}
\let\pscalculate\undefined
\usepackage{pst-bezier}
\usepackage{tikz,pgfplots,tkz-tab}
\usetikzlibrary{arrows,arrows.meta,patterns,calc,positioning,trees,decorations.markings,shapes.geometric,shapes.misc,shapes.arrows,intersections,decorations.pathmorphing,decorations.pathreplacing,backgrounds,shadows,babel}
\pgfplotsset{compat=1.18}
\usepgfplotslibrary{fillbetween}
\usepackage[np]{numprint}
\usepackage{listings,makecell,pas-tableur,twoopt}
\let\barre\undefined
\usepackage[pstricks]{tablvar}
\definecolor{darkgreen}{RGB}{0,100,0}
\providecommand{\red}{\color{red}}
\providecommand{\blue}{\color{blue}}
\providecommand{\green}{\color{green}}
\providecommand{\black}{\color{black}}
\providecommand{\white}{\color{white}}
\providecommand{\gray}{\color{gray}}
\providecommand{\cyan}{\color{cyan}}
\providecommand{\magenta}{\color{magenta}}
\providecommand{\yellow}{\color{yellow}}
\providecommand{\violet}{\color{violet}}
\providecommand{\orange}{\color{orange}}
\providecommand{\euro}{\texteuro}
\providecommand{\EUR}{\texteuro}
\providecommand{\ieme}{\textsuperscript{e}}
\providecommand{\degres}{\ensuremath{^\circ}}
\providecommand{\degree}{\ensuremath{^\circ}}
\providecommand{\vv}[1]{\overrightarrow{#1}}
\providecommand{\Coord}[1]{#1}
\providecommand{\decosix}{}
\providecommand{\starredbullet}{\ensuremath{\bullet}}
\pagestyle{empty}
\setlength{\parindent}{0pt}
\usepackage{adjustbox}\DeclareUnicodeCharacter{2236}{:}\DeclareUnicodeCharacter{2190}{\ensuremath{\leftarrow}}
\providecommand{\R}{\mathbb{R}}
\providecommand{\N}{\mathbb{N}}
\providecommand{\Z}{\mathbb{Z}}
\providecommand{\Q}{\mathbb{Q}}
\providecommand{\D}{\mathbb{D}}
\providecommand{\C}{\mathbb{C}}
\providecommand{\e}{{\mathrm{e}}}
\providecommand{\vect}[1]{\overrightarrow{#1}}
\providecommand{\vectt}[1]{\overrightarrow{\mathrm{#1}}}
\providecommand{\vvt}[1]{\overrightarrow{\mathrm{#1}}}
\providecommand{\np}[2][]{{#2~#1}}
\providecommand{\vv}[1]{{\overrightarrow{#1}}}
\renewcommand{\d}{{\mathrm{d}}}
\providecommand{\degres}{{^\circ}}
\providecommand{\degre}{{^\circ}}
\providecommand{\blue}{\color{blue}}
\providecommand{\red}{\color{red}}
\providecommand{\enskip}{\quad}
\providecommand{\llbracket}{\lbrack\!\lbrack}
\providecommand{\rrbracket}{\rbrack\!\rbrack}
\providecommand{\euro}{{\text{€}}}
\providecommand{\no}{{\text{n°}}}
\providecommand{\interval}[2]{{[#1;#2]}}
\providecommand{\intervFO}[2]{{[#1;#2[}}
\providecommand{\dotfill}{\ldots}

\providecommand{\ds}{\displaystyle}
\providecommand{\textding}[1]{#1}
\providecommand{\Oij}{\(\mathrm{O};\vec{\imath},\vec{\jmath}\)}
\providecommand{\Oijk}{\(\mathrm{O};\vec{\imath},\vec{\jmath},\vec{k}\)}
\providecommand{\Ouv}{\(\mathrm{O};\vec{u},\vec{v}\)}
\providecommand{\intervFF}[2]{\ensuremath{[#1;#2]}}
\renewcommand{\euro}{\ifmmode\text{\texteuro}\else\texteuro\fi}
\DeclareUnicodeCharacter{00B0}{\ifmmode^{\circ}\else\textdegree\fi}
\DeclareUnicodeCharacter{20AC}{\ifmmode\text{\texteuro}\else\texteuro\fi}
% Énoncés : Éducation nationale. Transcription et corrigés : APMEP, Denis Vergès et contributeurs. Découpage et classement : CoopMaths, Sébastien Lozano, Éric Elter et contributeurs. Adaptation de la présentation : Maths Explorer.
% CC BY-SA : https://coopmaths.fr/www/about/
% https://coopmaths.fr/alea/static/eam/2026/tex/eam_technologique_2026_06_polynesie_automatismes_cor.tex
\begin{document}
\noindent\textit{APMEP / CoopMaths. CC BY-SA. Présentation adaptée par Maths Explorer.}\par\medskip
\subsection*{Question 1}
	La règle est :\quad $\left(a^b\right)^c = a^{b \times c}$,\quad donc \quad $\left(5^{3}\right)^{2}=  5^{3\times 2} = 5^6$.

	Bonne réponse : \textbf{B.}\quad $5^6$.

\subsection*{Question 2}
	$ 1 + 25\% = 1+\dfrac{25}{100} = 1,25$.

	Bonne réponse : \textbf{C.}\quad 1,25.

\subsection*{Question 3}
	La droite n'est pas parallèle à l'axe des ordonnées. L'équation réduite est donc de la forme $y=mx+p$, où $p$ et $m$ sont à déterminer.

	Pour $x=0$, $y=0$ :\quad (la droite passe pour l'origine du repère) ainsi $p=0$.

	Pour $x=1$, $y=3$ :\quad ainsi $m$ est solution de $3 = m\times 1$, donc $m=3$:

	Bonne réponse : \textbf{C.}\quad $y=3x$.

\subsection*{Question 4}
	$\dfrac{1}{4} = \dfrac{25}{100} = 25~\%$. 25~\% des élèves \textbf{sont} adhérents à l'association sportive.

	Donc 75~\% des élèves ne sont pas adhérents à l'association sportive.

	Bonne réponse : \textbf{B.}\quad 75~\%.

\subsection*{Question 5}
	$\left(1+\dfrac{1}{2}\right)\times\dfrac{2}{3} = \left(\dfrac{2}{2}+\dfrac{1}{2}\right)\times\dfrac{2}{3} = \dfrac{3}{2}\times \dfrac{2}{3} = 1 $.

	Bonne réponse : \textbf{A.}\quad 1.

\subsection*{Question 6}
	On utilise l'identité remarquable :\quad $(2x+3)^2 = (2x)^2+2\times 2x \times 3 + 3^2 = 4x^2 + 12x + 9$.

	Bonne réponse : \textbf{B.}\quad $4x^2+12x+9$.


\subsection*{Question 7}

L'univers est $\Omega = \lbrace 1 ; 2 ; 3 ; 4 ; 5 ; 6 \rbrace$. L'évènement «~obtenir un numéro inférieur ou égal à 2~» est : $E = \lbrace 1 ; 2  \rbrace$.

Avec l'hypothèse d'équiprobabilité, $P(E) = \dfrac{2}{6} = \dfrac{1}{3}$.

Bonne réponse : \textbf{D.}\quad $ \dfrac{1}{3} $.

\subsection*{Question 8}

La moyenne se calcule comme le quotient de la somme des valeurs par le nombre de valeurs.

Bonne réponse : \textbf{A.}\quad $\dfrac{7+3+4+2+x}{5}=5$

\needspace{10\baselineskip}

\end{document}
